\chapter{Momentum}\label{ch:s1:momentum}

De mars à mai 2009, le facteur momentum américain a perdu 49.4~\%, et 52.8~\% sur l'année civile. Il était vendeur des actions qui avaient le plus chuté pendant la crise, et
ces actions ont rebondi le plus fort lorsque le marché s'est retourné. Le bilan à long terme du momentum avait fait figure de l'une des meilleures primes sur actions~: 7.4~\% par an depuis
1927, pour un ratio de Sharpe de 0.45. Mais ce bilan s'accompagne d'une asymétrie de $-3$, et en 1932 le facteur avait perdu 64.5~\% en une année. Normalisé chaque mois par sa
propre volatilité récente, le même facteur aurait perdu 16.8~\% en 2009 et doublé son ratio de Sharpe à long terme. Ce chapitre traite du momentum des actions, de ses
variantes et, surtout, de ses krachs. La construction est \texttt{firm.momstrat}.

\section{Momentum transversal}

\begin{definition}[Momentum transversal]\label{def:s1:momentum:cross}
Le \emph{momentum transversal}\index{momentum transversal} est la tendance des actions dont les rendements ont été élevés au cours des trois à douze derniers mois, par rapport
aux autres actions, à continuer de surperformer dans les mois suivants~; la stratégie standard classe les actions selon leur rendement entre douze mois et un mois en arrière (en écartant
le mois le plus récent, qui s'inverse) et achète les premières et vend les dernières.
\end{definition}

Jegadeesh et Titman l'ont documenté~: acheter les gagnantes passées et vendre les perdantes passées a rapporté des rendements significatifs sur des périodes de détention de trois à douze
mois, non expliqués par le risque systématique, et une partie du gain de la première année s'est estompée au cours des deux suivantes. Le mois écarté est le renversement du chapitre~2~: sur
un mois, les rendements s'inversent, et un signal de momentum qui l'inclut opère contre lui-même. Le facteur momentum de Kenneth French est l'indice de référence public~: les
portefeuilles à rendement passé élevé moins faible, parmi six portefeuilles pondérés par la valeur et répartis selon la taille et le rendement passé, reconstruits chaque mois.

\begin{center}
\small
\begin{tabular}{@{}lcc@{}}
\toprule
Facteur momentum français, 1927--2026 & brut & normalisé\\
\midrule
rendement annuel~; volatilité & 7.4~\%~; 16.3~\% & 15.4~\%~; 17.7~\%\\
ratio de Sharpe & 0.45 & 0.87\\
asymétrie~; kurtosis en excès (mensuel) & $-3.02$~; 27.3 & $-0.40$~; 3.1\\
pire mois & $-52.6~\%$ (8/1932) & $-29.1~\%$ (9/1939)\\
drawdown maximal & $-78.4~\%$ (6/1932 à 9/1939) & $-42.5~\%$ (6/1932 à 5/1933)\\
mars à mai 2009~; année civile 2009 & $-49.4~\%$~; $-52.8~\%$ & $-16.5~\%$~; $-16.8~\%$\\
\bottomrule
\end{tabular}
\end{center}

Sur le marché synthétique, qui introduit une dérive persistante dans chaque action (livre~7, chapitre~5), le signal 12--1 présente un coefficient d'information de rang mensuel de
0.020, et un portefeuille long--short par déciles (exposition brute de un, reconstruit tous les 21 jours) rapporte 4.1~\% par an net de dix points de base par unité échangée, avec un ratio de
Sharpe de 0.39 avant coûts et de 0.31 après, et un drawdown maximal de 40~\%. Il traite 5.7 fois son exposition brute par an, à sens unique.

\section{Momentum résiduel et momentum sectoriel}

\begin{definition}[Momentum résiduel, momentum sectoriel]\label{def:s1:momentum:residual}
Le \emph{momentum résiduel}\index{momentum résiduel} classe les actions selon la somme de leurs rendements résiduels du modèle factoriel sur la fenêtre de formation, divisée par
l'écart type des résidus, au lieu de les classer selon leurs rendements totaux. Le \emph{momentum sectoriel}\index{momentum sectoriel} classe les secteurs selon leurs rendements passés
et achète les actions des secteurs gagnants et vend celles des secteurs perdants.
\end{definition}

Le rendement sur douze mois d'une action s'explique surtout par ses facteurs~: une gagnante d'une année où les petites valeurs décotées ont progressé est une petite valeur décotée. Le momentum
sur rendements totaux porte donc les paris factoriels de l'année passée, et ceux-ci changent chaque mois. Blitz, Huij et Martens ont constaté que classer sur les rendements résiduels supprimait
ces expositions variables dans le temps et doublait à peu près les profits ajustés du risque, qui étaient aussi plus réguliers et moins concentrés sur les actions extrêmes. Moskowitz et
Grinblatt ont soutenu l'inverse en affirmant que les secteurs comptent~: un fort effet de momentum dans les composantes sectorielles des rendements expliquait une grande part du momentum
des actions individuelles, et les stratégies de \omterm{def:s1:momentum:residual}{momentum sectoriel} étaient très rentables à elles seules.

\begin{center}
\small
\begin{tabular}{@{}lccc@{}}
\toprule
marché synthétique, années 3 à 10 & rendement total & résiduel & sectoriel\\
\midrule
coefficient d'information de rang, mois suivant & 0.020 & 0.023 & 0.001\\
ratio de Sharpe avant, après 10 pb & 0.39, 0.31 & 0.55, 0.45 & $-0.04$, $-0.23$\\
rendement annuel après coûts~; volatilité & 4.1~\%~; 13.4~\% & 4.9~\%~; 11.0~\% & $-1.4~\%$~; 6.4~\%\\
drawdown maximal & $-40~\%$ & $-33~\%$ & $-23~\%$\\
\bottomrule
\end{tabular}
\end{center}

La version résiduelle l'emporte pour la raison avancée par Blitz et ses coauteurs~: la même dérive introduite, moins de bruit factoriel, une volatilité de 11.0~\% contre 13.4~\%. Le \omterm{def:s1:momentum:residual}{momentum
sectoriel} ne rapporte rien, car le marché synthétique n'en introduit aucun~: ses facteurs sectoriels sont indépendants d'un jour à l'autre. C'est la limite du marché synthétique,
et non une preuve contre Moskowitz et Grinblatt~; c'est aussi un rappel qu'un signal ne peut trouver que ce que son marché contient.

\begin{omfigure}
\begin{tikzpicture}
\begin{axis}[width=10cm,height=5.0cm,xlabel={\small année},ylabel={\small rendement net cumulé (\%)},tick label style={font=\footnotesize},
  xmin=2,xmax=10,grid=major,grid style={gray!25},legend pos=north west,legend style={font=\scriptsize},table/col sep=comma]
\addplot[omDef!60,thick] table[x=year,y=total]{figdata/strategies-1/05-momentum/books.csv};
\addplot[omDef,very thick] table[x=year,y=residual]{figdata/strategies-1/05-momentum/books.csv};
\addplot[omThm,thick] table[x=year,y=industry]{figdata/strategies-1/05-momentum/books.csv};
\legend{momentum à rendement total,momentum résiduel,momentum sectoriel}
\end{axis}
\end{tikzpicture}

\omcaption{Les trois portefeuilles de momentum sur le marché synthétique, net de dix points de base par unité échangée~: rendements cumulés (exposition brute de un). Données~:
\texttt{s1\_momentum.book}.}\label{fig:s1:momentum:books}
\end{omfigure}

\section{Momentum intrajournalier et momentum de l'heure du jour}

Le momentum opère aussi à l'échelle de la journée, sur le marché plutôt qu'entre les actions. Gao, Han, Li et Zhou ont établi, sur les données de l'ETF S\&P 500 de 1993 à 2013, que
le rendement de la première demi-heure du marché, mesuré depuis la clôture précédente, prédit son rendement de la dernière demi-heure, plus fortement les jours volatils, à fort volume, de récession et
d'annonces macroéconomiques, et dans dix autres ETF activement négociés. L'explication qu'ils privilégient combine un rééquilibrage peu fréquent des portefeuilles et des transactions tardives
d'investisseurs informés, proches de la clôture. Il s'agit d'une stratégie sur contrats à terme et ETF fondée sur des données intrajournalières, que le marché quotidien synthétique ne peut pas tester~; sa
fiche de stratégie repose sur le bilan publié, et le chapitre~25 revient sur les stratégies sur contrats à terme à court horizon.

\section{Krachs et leur gestion}

\begin{definition}[Krach du momentum]\label{def:s1:momentum:crash}
Un \emph{krach du momentum}\index{krach du momentum} est une courte période de très fortes pertes d'une stratégie de momentum, typiquement lorsque le marché rebondit vivement après
une baisse, car les perdantes passées que la stratégie vend à découvert (souvent des actions à bêta élevé fortement déprimées par la baisse) remontent davantage que les gagnantes passées qu'elle détient.
\end{definition}

Daniel et Moskowitz ont caractérisé les krachs~: des suites de rendements négatifs rares et persistantes, en partie prévisibles, survenant dans des états de panique,
après des baisses de marché et lorsque la volatilité est élevée, au moment même où le marché rebondit. Le mécanisme tient à un bêta que la stratégie acquiert sans
le choisir. Après une longue baisse, les perdantes passées sont les actions qui ont le plus chuté, c'est-à-dire celles au bêta le plus élevé~; un portefeuille de momentum est alors vendeur de bêta élevé,
et un rebond brutal le frappe de deux côtés. Les dix pires mois du facteur français sont six mois de 1931 à 1939, avril 2009, janvier 2001, novembre 2002 et
janvier 2023. Le marché synthétique compte trop peu de marchés baissiers pour montrer le mécanisme (le marché a baissé sur les deux années précédentes lors de seulement 6.7~\% de ses jours),
ce qui explique que cette section s'appuie sur les données françaises.

\begin{definition}[Momentum normalisé par la volatilité]\label{def:s1:momentum:scaled}
Le \emph{momentum normalisé par la volatilité}\index{momentum normalisé par la volatilité} (volatility-scaled momentum) détient un portefeuille de momentum dont la pondération est inversement proportionnelle à la volatilité réalisée récente de la stratégie elle-même,
de sorte que sa volatilité prévue soit constante~; la pondération est recalculée à chaque période à partir des données antérieures.
\end{definition}

Barroso et Santa-Clara ont montré que le risque du momentum est très variable et prévisible, et que le piloter élimine pratiquement les krachs et presque
double le ratio de Sharpe. La version du chapitre détient chaque mois une pondération de 12~\% divisée par la volatilité annualisée du facteur quotidien sur les 126
jours de bourse. Elle fait passer le ratio de Sharpe de 0.45 à 0.87, l'asymétrie de $-3.02$ à $-0.40$, et la perte de 2009 de 52.8~\% à 16.8~\%. Les krachs
surviennent après des mois volatils~: en avril 2009, lorsque le facteur brut a perdu 34.4~\%, la pondération normalisée valait 0.29, et la perte normalisée 10.1~\%. La normalisation n'aide pas
lorsque le krach vient d'une période calme~: en septembre 1939, la pondération valait 0.92, et le facteur normalisé a connu son pire mois ($-29.1~\%$). Sur le marché synthétique,
où la volatilité du momentum varie moins, la normalisation apporte peu~: le ratio de Sharpe du portefeuille en rendement total reste à 0.52 sur l'échantillon normalisé et son
drawdown maximal passe de 30~\% à 25~\%.

\begin{omfigure}
\begin{tikzpicture}
\begin{axis}[width=12.5cm,height=5.4cm,xlabel={\small année},ylabel={\small rendement de l'année civile (\%)},tick label style={font=\footnotesize},
  xmin=1927,xmax=2026,ymin=-70,ymax=150,grid=major,grid style={gray!25},legend pos=north west,legend style={font=\scriptsize},table/col sep=comma,
  x tick label style={/pgf/number format/1000 sep={}}]
\addplot[omThm,thick,mark=*,mark size=0.9pt] table[x=year,y=raw_pct]{figdata/strategies-1/05-momentum/annual.csv};
\addplot[omDef,thick,mark=square*,mark size=0.9pt] table[x=year,y=scaled_pct]{figdata/strategies-1/05-momentum/annual.csv};
\legend{brut,normalisé par la volatilité}
\end{axis}
\end{tikzpicture}

\omcaption{Rendements annuels civils du facteur momentum de Kenneth French, brut et normalisé chaque mois à 12~\% par sa volatilité glissante sur 126 jours, 1928--2025.
Dérivé de la Kenneth R. French Data Library (Mom, mensuel et quotidien, fichier CRSP 202607)~; la série brute n'est pas
redistribuée.}\label{fig:s1:momentum:annual}
\end{omfigure}

La version normalisée n'est pas gratuite. Elle recourt à l'effet de levier chaque fois que la volatilité glissante du momentum est inférieure à 12~\% (une grande partie de son rendement moyen de 15.4~\%
contre 7.4~\% pour le facteur brut provient de ce levier), et un fonds qui la met en œuvre a besoin de la capacité et du financement correspondants. Or l'amélioration est mesurée sur le siècle même qui l'a suggérée~; Daniel et
Moskowitz, dans leur version dynamique qui prévoit aussi la moyenne du momentum à partir de l'état du marché, rapportent un doublement comparable sur plusieurs marchés et classes d'actifs.

\section{Mise en œuvre~: taux de rotation, saut d'un mois, coûts}

Le momentum est une stratégie lente selon les standards de cette partie du livre, et ses coûts sont faibles par unité de capital~: le portefeuille synthétique par déciles négocie 5.7 fois
son volume brut annuel à sens unique, ce qui, à dix points de base par unité négociée, coûte 1.14~\% par an, contre 8~\% à 38~\% pour les portefeuilles quotidiens des chapitres~2 et~3. Trois
choix de mise en œuvre comptent davantage que le taux de frais. Le mois sauté supprime le retournement, comme plus haut. La fréquence de rééquilibrage arbitre entre la fraîcheur
contre la rotation (le rythme mensuel est la convention). Et la pondération (égale ou par capitalisation, déciles ou terciles) déplace le portefeuille vers des petites valeurs, moins liquides,
où le momentum est plus fort et la négociation plus chère~; le facteur français moyenne grandes et petites valeurs pour les garder toutes deux.

\section{Fiches de stratégie}

\begin{strategyfile}{Momentum transversal 12-1}\label{strat:s1:momentum:classic}
\sfield{Qui paie, et pourquoi} Des investisseurs qui sous-réagissent aux nouvelles puis les poursuivent~; le portefeuille est en avance sur la tendance et en retard sur sa fin.
\sfield{Instruments et plateformes} Actions liquides, achat des gagnantes et vente à découvert des perdantes, ou biais à l'achat seul.
\sfield{Signal} Rendement entre il y a douze mois et il y a un mois, classé sur l'ensemble des actions.
\sfield{Taille des positions et exécution} Déciles ou terciles, reconstruits chaque mois~; pondérations par capitalisation ou égales.
\sfield{Coûts} Plusieurs fois la valeur du portefeuille négociée par an~; faibles par unité de capital.
\sfield{Comment elle s'éteint} Krachs après les marchés baissiers lorsque les perdantes rebondissent (2009, 1932)~; érosion à mesure que les capitaux s'y massent (comomentum, livre~7, chapitre~28).
\sfield{Horizon, capacité, infrastructure} Mois~; grande capacité~; un classement mensuel.
\sfield{Backtest rigoureux} Univers à chaque date, y compris les actions radiées~; le mois sauté~; les années de krach dans l'échantillon.
\sfield{Sources} Jegadeesh et Titman (1993)~; facteur momentum de Kenneth French, 1927--2026~: 7.4~\% par an, ratio de Sharpe de 0.45 avant coûts, drawdown maximal de 78.4~\%.
\end{strategyfile}

\begin{strategyfile}{Momentum résiduel}\label{strat:s1:momentum:residual}
\sfield{Qui paie, et pourquoi} Comme pour le momentum, avec la part propre à chaque action de la tendance isolée.
\sfield{Instruments et plateformes} Comme pour le momentum.
\sfield{Signal} Somme des rendements résiduels du modèle factoriel sur les mois $-12$ à $-2$, divisée par leur écart-type.
\sfield{Taille des positions et exécution} Comme pour le momentum~; le portefeuille est, par construction, plus proche de la neutralité aux facteurs.
\sfield{Coûts} Comme pour le momentum.
\sfield{Comment elle s'éteint} Comme le momentum, avec des krachs plus petits puisque le portefeuille porte moins de bêta.
\sfield{Horizon, capacité, infrastructure} Mois~; un modèle factoriel.
\sfield{Backtest rigoureux} Résidus d'un modèle estimé uniquement sur des données passées.
\sfield{Sources} Blitz, Huij et Martens (2011)~: profits ajustés du risque environ le double de ceux du momentum en rendement total.
\end{strategyfile}

\begin{strategyfile}{Momentum sectoriel}\label{strat:s1:momentum:industry}
\sfield{Qui paie, et pourquoi} Diffusion lente de l'information à l'échelle des secteurs.
\sfield{Instruments et plateformes} Paniers sectoriels ou fonds indiciels cotés (ETF), ou actions par secteur.
\sfield{Signal} Rendements passés des secteurs sur six à douze mois.
\sfield{Taille des positions et exécution} Achat des secteurs en tête, vente à découvert de ceux en queue~; mensuel.
\sfield{Coûts} Faibles lorsqu'elle est négociée via des fonds sectoriels.
\sfield{Comment elle s'éteint} Comme le momentum~; paris concentrés sur peu de secteurs.
\sfield{Horizon, capacité, infrastructure} Mois~; grande capacité.
\sfield{Backtest rigoureux} Classifications sectorielles telles qu'elles étaient à chaque date.
\sfield{Sources} Moskowitz et Grinblatt (1999).
\end{strategyfile}

\begin{strategyfile}{Momentum infrajournalier de la dernière demi-heure}\label{strat:s1:momentum:intraday}
\sfield{Qui paie, et pourquoi} Les investisseurs qui rééquilibrent en fin de séance et les traders tardivement informés qui négocient dans le sens du mouvement précoce de la journée.
\sfield{Instruments et plateformes} Contrats à terme sur indices et ETF.
\sfield{Signal} Rendement de la première demi-heure depuis la clôture précédente.
\sfield{Taille des positions et exécution} Position sur la dernière demi-heure dans son sens~; à plat à la clôture.
\sfield{Coûts} Un aller-retour par jour sur les instruments les plus liquides.
\sfield{Comment elle s'éteint} Concurrence dans la demi-heure de clôture~; changements dans l'identité de ceux qui rééquilibrent à la clôture.
\sfield{Horizon, capacité, infrastructure} Une demi-heure~; données infrajournalières et exécution.
\sfield{Backtest rigoureux} Exécuter aux cours disponibles au début de la dernière demi-heure, et non à leur moyenne.
\sfield{Sources} Gao, Han, Li et Zhou (2018)~: prévisibilité significative dans les données de l'ETF S\&P 500, 1993--2013.
\end{strategyfile}

\begin{strategyfile}{Momentum normalisé par la volatilité}\label{strat:s1:momentum:scaled}
\sfield{Qui paie, et pourquoi} Comme pour le momentum~; la normalisation retire un risque que la prime ne rémunère pas.
\sfield{Instruments et plateformes} Comme pour le momentum, avec un effet de levier disponible.
\sfield{Signal} La pondération du portefeuille de momentum~: une volatilité cible rapportée à sa volatilité réalisée glissante.
\sfield{Taille des positions et exécution} Recalculée chaque mois (ou plus souvent) à partir des rendements quotidiens antérieurs au mois.
\sfield{Coûts} Ceux du portefeuille de momentum, plus ceux liés aux variations de pondération.
\sfield{Comment elle s'éteint} Krachs issus d'états calmes (septembre 1939)~; contraintes liées au levier les années calmes.
\sfield{Horizon, capacité, infrastructure} Mois~; rendements quotidiens de la stratégie.
\sfield{Backtest rigoureux} Pondérations à partir des données antérieures à chaque mois~; financement du levier.
\sfield{Sources} Barroso et Santa-Clara (2015)~; Daniel et Moskowitz (2016)~; données de Kenneth French~: ratio de Sharpe de 0.45 à 0.87 et perte de 2009 de 16.8~\% au lieu de
52.8~\%, 1927--2026, avant coûts.
\end{strategyfile}

\section{Tutoriel~: printemps 2009}\label{tut:s1:momentum:tut}

\begin{tutorial}
\textbf{Objectif.} Construire trois portefeuilles de momentum sur le marché synthétique, puis mesurer le krach de 2009 dans le facteur français, brut et normalisé. \textbf{État final~:} les
deux tableaux et la \cref{fig:s1:momentum:annual}.
\begin{tutsteps}
\item \textbf{Signaux}~: \omterm{def:s1:momentum:residual}{momentum résiduel} et \omterm{def:s1:momentum:residual}{momentum sectoriel}.
\omcode{code/firm/momstrat/firm_momstrat.py}{45}{61}{Momentum résiduel et momentum sectoriel.}\label{lst:s1:momentum:signals}
\item \textbf{Portefeuilles}~: \texttt{s1\_momentum.book(kind)} et \texttt{summary(kind)} pour les trois signaux.
\item \textbf{Normalisation}~: une pondération issue de la volatilité glissante propre à la stratégie.
\omcode{code/firm/momstrat/firm_momstrat.py}{79}{84}{Normalisation par la volatilité.}\label{lst:s1:momentum:scale}
\item \textbf{Données de French}~: \texttt{s1\_fetch\_mom.py} une fois, puis \texttt{french()}~; \texttt{fig\_momentum.py}.
\end{tutsteps}
\textbf{Pour aller plus loin.} Conditionner la normalisation à l'état du marché (Daniel et Moskowitz)~; utiliser des terciles et des pondérations par capitalisation comme le facteur français~; ajouter
le retournement du chapitre~2 au mois sauté, et mesurer ce qu'il coûte.
\end{tutorial}

\section{Construction~: stratégies de momentum}\label{bld:s1:momentum:momstrat}

\begin{build}
\bfield{Rôle} La famille des signaux de momentum et la gestion des krachs dont tout portefeuille de momentum a besoin.
\bfield{Interface} \texttt{total\_mom(ret, lookback, skip)}, \texttt{residual\_mom(E, lookback, skip)},
\texttt{industry\_mom(ret, industry, cap, lookback, skip)}, \texttt{decile\_book(signal, universe, q)}, \texttt{vol\_scale(r, target, window, periods)}, \texttt{bear(mkt, window)}.
\bfield{Règles} Signaux et pondérations issus des données antérieures à chaque transaction~; le mois sauté~; normalisation à partir des rendements propres à la stratégie.
\bfield{Tests d'acceptation} \texttt{code/firm/momstrat/tests/}~: la fenêtre de saut vérifiée à la main~; signaux résiduels et sectoriels sur des données plantées~; pondérations par déciles~; la
pondération de la normalisation à la main~; l'indicateur d'état du marché.
\bfield{Pour aller plus loin} Pondérations dynamiques issues de prévisions de la moyenne et de la variance (Daniel et Moskowitz)~; momentum temporel (time-series momentum) sur contrats à terme (chapitre~19).
\end{build}

\begin{omsources}
\item N.~Jegadeesh et S.~Titman, «~Returns to buying winners and selling losers~», \emph{Journal of Finance} 48(1), 1993.
\item K.~Daniel et T.~J.~Moskowitz, «~Momentum crashes~», \emph{Journal of Financial Economics} 122(2), 2016.
\item P.~Barroso et P.~Santa-Clara, «~Momentum has its moments~», \emph{Journal of Financial Economics} 116(1), 2015.
\item D.~Blitz, J.~Huij et M.~Martens, «~Residual momentum~», \emph{Journal of Empirical Finance} 18(3), 2011.
\item T.~J.~Moskowitz et M.~Grinblatt, «~Do industries explain momentum~?~», \emph{Journal of Finance} 54(4), 1999.
\item L.~Gao, Y.~Han, S.~Z.~Li et G.~Zhou, «~Market intraday momentum~», \emph{Journal of Financial Economics} 129(2), 2018.
\item Kenneth R. French Data Library, Momentum Factor (Mom), données mensuelles et quotidiennes.
\end{omsources}

\section{Exercices}

\begin{exercise}[$\star$]\label{exo:s1:momentum:1}
Calculer la pondération que la normalisation du chapitre attribue au facteur de momentum, sachant que sa volatilité glissante est de 40~\% par an, et ce que deviendrait une perte brute de 34~\%.
\end{exercise}

\begin{exercise}[$\star$]\label{exo:s1:momentum:2}
Calculer le rendement qui compense une perte de 50~\%. Puis, après une perte de 50~\% suivie d'un gain de 100~\%, calculer où se situe la stratégie.
\end{exercise}

\begin{exercise}[$\star$]\label{exo:s1:momentum:3}
Calculer le coût annuel, pour le portefeuille par déciles, de dix points de base par unité négociée, sachant qu'il négocie chaque année, à sens unique, 5.7 fois son exposition brute.
\end{exercise}

\begin{exercise}[$\star\star$]\label{exo:s1:momentum:4}
Expliquer pourquoi un portefeuille de momentum est vendeur de titres à bêta élevé après une longue baisse du marché, et ce qui se passe lorsque le marché rebondit.
\end{exercise}

\begin{exercise}[$\star\star$]\label{exo:s1:momentum:5}
Expliquer pourquoi le momentum écarte le mois le plus récent.
\end{exercise}

\begin{exercise}[$\star\star$]\label{exo:s1:momentum:6}
Commenter ce que la normalisation par la volatilité peut et ne peut pas faire, à partir du pire mois du facteur normalisé (septembre 1939), survenu avec une pondération de 0.92.
\end{exercise}

\begin{exercise}[$\star\star\star$]\label{exo:s1:momentum:7}
\emph{Programmation.} Exécuter \texttt{scaled('residual')}. Comparer les ratios de Sharpe avec celui du portefeuille à rendement total, et expliquer pourquoi la normalisation aide le facteur français
davantage que les portefeuilles synthétiques.
\end{exercise}

\begin{exercise}[$\star\star\star$]\label{exo:s1:momentum:8}
\emph{Trouver la faille.} «~Le momentum a rapporté 7~\% par an pendant un siècle avec un ratio de Sharpe de 0.45~; à trois fois l'effet de levier, c'est une stratégie à 21~\%.~»
\end{exercise}

\section{Problème~: printemps 2009}

\begin{problem}[{Problème du week-end --- le krach et son remède}]\label{pb:s1:momentum:1}
Le facteur de momentum français, brut et normalisé, et trois portefeuilles de momentum sur \texttt{firm.synthmkt}.

\medskip\noindent\textbf{Partie I --- La prime.}
\begin{enumerate}
\item Définir le \omterm{def:s1:momentum:cross}{momentum transversal} et le signal 12--1.
\item Indiquer ce que Jegadeesh et Titman ont trouvé.
\item Donner le rendement, la volatilité, le ratio de Sharpe et l'asymétrie du facteur français depuis 1927.
\item Indiquer son pire mois et son drawdown maximal.
\end{enumerate}

\medskip\noindent\textbf{Partie II --- Les variantes.}
\begin{enumerate}[resume]
\item Définir le \omterm{def:s1:momentum:residual}{momentum résiduel} et le \omterm{def:s1:momentum:residual}{momentum sectoriel}, et donner leurs résultats publiés.
\item Donner les IC et les ratios de Sharpe des trois portefeuilles synthétiques.
\item Expliquer pourquoi le \omterm{def:s1:momentum:residual}{momentum résiduel} l'emporte ici.
\item Expliquer pourquoi le \omterm{def:s1:momentum:residual}{momentum sectoriel} ne rapporte rien ici.
\end{enumerate}

\medskip\noindent\textbf{Partie III --- Krachs.}
\begin{enumerate}[resume]
\item Définir un \omterm{def:s1:momentum:crash}{krach du momentum} et décrire quand les krachs se produisent.
\item Expliquer le mécanisme en termes de bêta.
\item Indiquer quels ont été les dix pires mois.
\item Expliquer pourquoi le marché synthétique ne reproduit pas le mécanisme.
\end{enumerate}

\medskip\noindent\textbf{Partie IV --- Le verdict.}
\begin{enumerate}[resume]
\item Énoncer le \emph{résultat nommé}~: le drawdown du facteur de momentum en 2009 dans les données de French et celui de la version normalisée par la volatilité.
\item Indiquer l'effet de la normalisation sur le ratio de Sharpe de long terme et sur l'asymétrie.
\item Indiquer quand la normalisation échoue.
\item Indiquer le niveau d'effet de levier qu'exige la version normalisée.
\item Indiquer ce que Gao et ses coauteurs ont trouvé à l'intérieur de la journée.
\item Indiquer ce qu'il faudrait surveiller sur un portefeuille de momentum en production.
\item Évaluer l'importance des coûts du momentum par rapport à ceux du chapitre~2.
\item Donner en une phrase le risque du momentum.
\end{enumerate}
\end{problem}

\section{Questions d'entretien}

\begin{interviewq}[$\star$ \iqroles{chercheur, trader}]\label{iq:s1:momentum:1}
Qu'est-ce que le momentum, et pourquoi saute-t-on le dernier mois~?
\end{interviewq}

\begin{interviewq}[$\star\star$ \iqroles{chercheur}]\label{iq:s1:momentum:2}
Pourquoi le momentum subit-il des krachs, et quand~?
\end{interviewq}

\begin{interviewq}[$\star\star$ \iqroles{chercheur, risque}]\label{iq:s1:momentum:3}
Comment géreriez-vous le risque de \omterm{def:s1:momentum:crash}{krach du momentum}~?
\end{interviewq}

\begin{interviewq}[$\star\star$ \iqroles{chercheur}]\label{iq:s1:momentum:4}
\omterm{def:s1:momentum:residual}{Momentum résiduel} par rapport au momentum à rendement total~: quelle est la différence, et lequel négocieriez-vous~?
\end{interviewq}

\begin{interviewq}[$\star\star$ \iqroles{risque}]\label{iq:s1:momentum:5}
Le bêta de votre portefeuille de momentum était de 0.1 il y a un mois et est de $-0.4$ aujourd'hui. Qu'est-il arrivé, et que faites-vous~?
\end{interviewq}

\begin{interviewq}[$\star\star\star$ \iqroles{chercheur}]\label{iq:s1:momentum:6}
Montrer que normaliser une stratégie par l'inverse d'une prévision de sa volatilité augmente son ratio de Sharpe lorsque la volatilité est prévisible et sans lien avec la moyenne.
Dans quel cas peut-elle réduire le ratio~?
\end{interviewq}
